%======================================================================
% 附录K：跨教材经典例题精选（注明出处）
%======================================================================
% 章标题 \appchapter{附录K：...} 由主文件发出，本文件只提供正文。

本附录从 Chen、Fitzpatrick、Goldston \& Rutherford 等经典教材中挑八道\textbf{最能检验基础图像}的题，逐题给出完整解答与出处。

选题只看三条：只用本书第 1--12 章的工具，答案能\textbf{独立复算}，而且背后都有一个「判断型」的物理结论，不只是代数练习。

\textbf{使用建议：}每道题先合书独立做一遍（动作五），再看解答。八道题覆盖八个不同的物理环节，适合按附录 I 的「交错练习」方式打乱使用。

\begin{tip}{复算与自查}{tip-appk-recalc}
若你的结果与文中差一个量级，几乎总是以下几个原因之一：用错了温度（电子 / 离子、eV / K）、漏掉平方根、或把 SI 与 CGS 混用。建议对照附录 F 的高频易错点清单一起看。
\end{tip}


\section{德拜屏蔽与等离子体参数}

\begin{example}{德拜长度、$N_D$ 与屏蔽势的衰减}{}
\textbf{题目：}某氢等离子体 $T_e=10\ \mathrm{eV}$、$n=10^{18}\ \mathrm{m^{-3}}$。
（a）求德拜长度 $\lambda_D$。
（b）求德拜球内的粒子数 $N_D$，并判断能否用流体（连续介质）图像描述。
（c）把测试电荷放在原点，$r=5\lambda_D$ 处的屏蔽势是裸库仑势的百分之几？

\textbf{解：}
（a）$\lambda_D=\sqrt{\dfrac{\eps_0 k_BT_e}{ne^2}}=7430\sqrt{\dfrac{T_e[\mathrm{eV}]}{n[\mathrm{m^{-3}}]}}\ \mathrm{m}
=7430\sqrt{\dfrac{10}{10^{18}}}=2.35\times10^{-5}\ \mathrm{m}\approx23.5\ \mathrm{\mu m}$。

（b）$N_D=\dfrac{4\pi}{3}n\lambda_D^3
=\dfrac{4\pi}{3}\times10^{18}\times(2.3508\times10^{-5})^3\approx5.4\times10^{4}$。
$N_D\gg1$，德拜球内粒子数极多，统计描述成立，流体图像在这里是合法的。

（c）屏蔽势 $\phi(r)=\dfrac{q}{4\pi\eps_0 r}\ee^{-r/\lambda_D}$，比值即 $\ee^{-5}=6.7\times10^{-3}$，只有 $0.67\%$。

{\footnotesize
\textbf{出处：}Chen《Introduction to Plasma Physics and Controlled Fusion》第 3 版第 1 章（德拜屏蔽）。
Fitzpatrick 第 2 版第 1 章 Debye Shielding。
}

\textbf{常见错误：}用离子温度算 $\lambda_D$（判据由电子决定）；漏掉 $N_D$ 中的 $4\pi/3$ 因子；把「$\lambda_D$ 内只有一个粒子」当成判据（正确判据是 $N_D\gg1$）。
\end{example}


\section{磁镜与损失锥}

\begin{example}{磁镜比 $R=10$ 的临界角与损失锥比例}{}
\textbf{题目：}磁镜的磁镜比 $R=B_{\max}/B_{\min}=10$。
（a）求临界投掷角 $\theta_c$。
（b）各向同性速度分布中，有多少比例的粒子落在损失锥内？
（c）若把 $R$ 提高到 $100$，该比例变为多少？

\textbf{解：}
（a）由 $\sin^2\theta_c=B_{\min}/B_{\max}=1/R$ 得。
$\theta_c=\arcsin\sqrt{1/10}=18.43^\circ$。

（b）各向同性分布中投掷角概率密度正比于 $\sin\theta$，故损失锥内的比例为。
$P=\displaystyle\int_0^{\theta_c}\sin\theta\,\diff\theta=1-\cos\theta_c
=1-\sqrt{1-\tfrac1R}=1-\sqrt{0.9}=5.13\%$。

（c）$R=100$ 时 $P=1-\sqrt{0.99}=0.50\%$：磁镜比提高十倍，损失锥比例降到约十分之一。

{\footnotesize
\textbf{出处：}Chen 第 2 章（单粒子运动，磁镜）。
Fitzpatrick 第 2 版第 2 章 Magnetic Mirrors。
}

\textbf{常见错误：}写成 $\sin\theta_c=1/R$（漏掉开方）；用 $v_\parallel/v_\perp$ 而不是 $\sin^2\theta$ 表达判据；把损失锥比例当成 $\theta_c/90^\circ$ 的线性比例。
\end{example}


\section{梯度漂移与曲率漂移的方向}

\begin{example}{$\grad B$ 漂移的方向判别与电荷分离}{}
\textbf{题目：}设磁场沿 $z$ 方向、大小沿 $x$ 线性减弱：$\vect{B}=B_0(1-x/L)\unitvec{z}$，$B_0>0$，$L>0$。
（a）求 $\grad B$，并判断离子与电子的 $\grad B$ 漂移方向。
（b）说明在托卡马克中，$\grad B$ 漂移与曲率漂移为什么会导致电荷分离，以及它是如何被抑制的。

\textbf{解：}
（a）$\grad B=-\dfrac{B_0}{L}\unitvec{x}$，于是
$\vect{B}\times\grad B=B_0\left(1-\dfrac{x}{L}\right)\unitvec{z}\times\left(-\dfrac{B_0}{L}\unitvec{x}\right)
=-\dfrac{B_0^2}{L}\left(1-\dfrac{x}{L}\right)\unitvec{y}$。
代入 $\vect{v}_{\grad B}=\dfrac{mv_\perp^2}{2qB}\dfrac{\vect{B}\times\grad B}{B^2}$：
离子（$q>0$）沿 $-\unitvec{y}$ 漂移，电子（$q<0$）沿 $+\unitvec{y}$ 漂移。方向相反，与电荷符号有关。

（b）在托卡马克中 $B\propto1/R$，$\grad B$ 指向内侧。这里最容易错的一步是\textbf{把曲率矢量与漂移方向混为一谈}：曲率矢量 $\kappa=(\unitvec{b}\cdot\grad)\unitvec{b}$ 确实指向曲率中心（环向场下 $\kappa=-\unitvec{R}/R_c$），但曲率漂移沿 $\unitvec{b}\times\kappa$，\emph{垂直于}曲率矢量：
\[
\vect{v}_R=\frac{mv_\parallel^2}{qB}\,(\unitvec{b}\times\kappa),
\qquad
\unitvec{b}\times\kappa=\unitvec{\varphi}\times\left(-\frac{\unitvec{R}}{R_c}\right)=+\frac{\unitvec{z}}{R_c}.
\]
漂移方向是竖直的（正离子向上、电子向下），\textbf{不是}指向曲率中心。
它与 $\grad B$ 漂移的\textbf{竖直分量同向}（后者也沿 $\pm\unitvec{z}$），且都因含 $1/q$ 而对电子、离子反向，于是等离子体上下两侧积累异号电荷，形成竖直电场；该电场与环向磁场产生 $\vect{E}\times\vect{B}$ 漂移，把粒子向外推。
抑制机制是\textbf{旋转变换}（磁力线的环向扭曲）：粒子沿磁力线运动时上下交替，短接了电荷分离。托卡马克因此需要足够大的等离子体电流，「磁通管」图像中才会出现 Pfirsch--Schl\"uter 电流。

{\footnotesize
\textbf{出处：}Chen 第 2 章（单粒子运动，非均匀磁场中的漂移）。
Fitzpatrick 第 2 版第 2 章 Magnetic Drifts。

托卡马克部分见 Wesson《Tokamaks》（按主题）。
}

\textbf{常见错误：}认为 $\grad B$ 漂移与电荷符号无关。$\vect{B}\times\grad B$ 确实与 $q$ 无关，但漂移速度中含 $1/q$，所以方向相反。另一种错法是把曲率漂移的方向当作曲率矢量的方向：漂移 $\propto\unitvec{b}\times\kappa$，与 $\kappa$ 垂直，两者不可互换。还有人把曲率漂移与 $\grad B$ 漂移在本例的几何下判成反向，其实它们同向；也有人以为电荷分离无法消除。
\end{example}


\section{Alfvén 速度与磁冻结}

\begin{example}{Alfvén 速度的量级与磁场扰动的传播时间}{}
\textbf{题目：}氘等离子体 $n=10^{20}\ \mathrm{m^{-3}}$、$B=5\ \mathrm{T}$。
（a）求 Alfvén 速度 $v_A$。（b）它是光速的百分之几？
（c）磁场扰动沿 $1\ \mathrm{m}$ 传播需要多久？（d）若密度增大到 $10^{22}\ \mathrm{m^{-3}}$，$v_A$ 变为多少？

\textbf{解：}
（a）质量密度由离子提供：$\rho\approx nm_D=10^{20}\times3.368\times10^{-27}=3.37\times10^{-7}\ \mathrm{kg\,m^{-3}}$。
$v_A=\dfrac{B}{\sqrt{\mu_0\rho}}
=\dfrac{5}{\sqrt{1.2566\times10^{-6}\times3.368\times10^{-7}}}
=7.69\times10^{6}\ \mathrm{m\,s^{-1}}$。

（b）$v_A/c=7.69\times10^{6}/2.998\times10^{8}=2.6\%$。这里只能得出一个结论：\textbf{非相对论 Alfvén 速度公式适用}（位移电流可以忽略），\textbf{不能据此断定磁冻结成立}。磁冻结要求理想欧姆定律 $\vect{E}+\vect{v}\times\vect{B}=0$ 成立，至少要磁雷诺数 $R_m=\mu_0Lv/\eta\gg1$，还要检查霍尔项与电子压强项是否可丢（第 5 章「MHD 适用条件：三道独立的门槛」）。本题没有给温度、电阻率与流动尺度，用现有输入无法判断：\emph{同一个} $v_A/c$ 在强碰撞的电阻性等离子体里成立，在理想 MHD 等离子体里也成立。作为对照：取 $L=1\ \mathrm{m}$、$\eta(10\ \mathrm{eV})=1.64\times10^{-5}\ \Omega\cdot\mathrm{m}$（K.6 已算），则
\[
R_m=\frac{\mu_0Lv_A}{\eta}
=\frac{1.2566\times10^{-6}\times1\times7.69\times10^{6}}{1.64\times10^{-5}}
\approx5.9\times10^{5}\gg1,
\]
此时才能谈磁冻结。而电阻扩散时标 $\tau_\eta=\mu_0L^2/\eta\approx0.076\ \mathrm{s}$ 给出了一个定量边界：在 $L\sim1\ \mathrm{m}$ 的尺度上，静止导电流体的磁场正是以这个时标扩散掉的（见 K.6）。

（c）$t=L/v_A=1/7.69\times10^{6}=1.30\times10^{-7}\ \mathrm{s}$。

（d）$v_A\propto1/\sqrt{\rho}\propto1/\sqrt{n}$：密度增大 100 倍，$v_A$ 减小到 $1/10$，即 $7.7\times10^{5}\ \mathrm{m\,s^{-1}}$。

{\footnotesize
\textbf{出处：}Fitzpatrick 第 2 版第 8 章（MHD 流体：磁压、磁通冻结、MHD 波）。
Chen 第 4 章（等离子体中的波）。
}

\textbf{常见错误：}用电子质量计算 $\rho$，而惯性由离子提供。忘记 $\mu_0$ 的平方根。把 $v_A$ 与离子声速 $c_s$ 混用：后者由电子温度决定，与 $B$ 无关。把 $v_A/c\ll1$ 当成磁冻结的充分条件，而它只保证非相对论公式适用；$R_m=\mu_0Lv/\eta\gg1$ 才是丢掉电阻项的条件。
\end{example}


\section{Landau 阻尼率}

\begin{example}{朗缪尔波的 Landau 阻尼率与每周期衰减}{}
\textbf{题目：}朗缪尔波满足 $k\lambda_D=0.3$。
（a）用 $\dfrac{\gamma}{\omega}=-\sqrt{\dfrac{\pi}{8}}\,(k\lambda_D)^{-3}
\exp\!\left(-\dfrac{1}{2k^2\lambda_D^2}-\dfrac{3}{2}\right)$ 求阻尼率。
（b）估算一个波周期内振幅衰减的百分比。
（c）为什么 $k\lambda_D\ll1$ 时阻尼可以忽略？

\textbf{解：}
（a）$(k\lambda_D)^{-3}=0.3^{-3}=37.04$，
$\exp[-1/(2\times0.09)]=\ee^{-5.556}=3.87\times10^{-3}$，$\ee^{-3/2}=0.2231$，
故 $\dfrac{\gamma}{\omega}=-0.6267\times37.04\times3.87\times10^{-3}\times0.2231=-2.00\times10^{-2}$。

（b）一个周期内振幅因子为 $\exp\!\left(2\pi\dfrac{\gamma}{\omega}\right)=\ee^{-0.126}=0.882$，即衰减约 $11.8\%$。

（c）指数因子 $\exp[-1/(2k^2\lambda_D^2)]$ 随 $k\lambda_D$ 减小而\textbf{双重指数}衰减：$k\lambda_D=0.1$ 时它已是 $\ee^{-50}\approx2\times10^{-22}$。物理上，长波相速度远大于热速度，能与波共振的粒子数指数稀少，因此无碰撞阻尼可以忽略。

{\footnotesize
\textbf{出处：}Fitzpatrick 第 2 版第 7 章 Landau Damping。
Chen 第 7 章（动理学理论）。
}

\textbf{常见错误：}把 $\gamma/\omega$ 直接当成「每周期衰减比例」（应乘以 $2\pi$）；把 $\ee^{-3/2}$ 漏掉；误认为 Landau 阻尼来自碰撞；它是无碰撞的相混合，不增加熵（在线性阶段）。
\end{example}


\section{Spitzer 电阻率与磁场扩散时间}

\begin{example}{从温度估算电阻率与磁场扩散时间}{}
\textbf{题目：}某托卡马克等离子体 $T=1\ \mathrm{keV}$、$Z=1$、$\ln\Lambda=17$，特征长度 $L=1\ \mathrm{m}$。
（a）用 $\eta=5.2\times10^{-5}\dfrac{Z\ln\Lambda}{T[\mathrm{eV}]^{3/2}}\ \Omega\cdot\mathrm{m}$ 求电阻率。
（b）求磁场扩散时间 $\tau_\eta=\mu_0L^2/\eta$。
（c）若温度降到 $10\ \mathrm{eV}$（$\ln\Lambda=10$），$\eta$ 与 $\tau_\eta$ 各变为多少？

\textbf{解：}
（a）$T^{3/2}=1000^{1.5}=3.16\times10^{4}$，
$\eta=5.2\times10^{-5}\times17/3.16\times10^{4}=2.80\times10^{-8}\ \Omega\cdot\mathrm{m}$。

（b）$\tau_\eta=\dfrac{1.2566\times10^{-6}\times1}{2.80\times10^{-8}}=44.9\ \mathrm{s}$。

（c）$10^{1.5}=31.6$，$\eta=5.2\times10^{-5}\times10/31.6=1.64\times10^{-5}\ \Omega\cdot\mathrm{m}$（约为原来的 $590$ 倍），
$\tau_\eta=1.2566\times10^{-6}/1.64\times10^{-5}=7.6\times10^{-2}\ \mathrm{s}$。

\textbf{结论：}温度提高两个数量级，$\eta$ 下降约 $590$ 倍，磁场扩散时间相应延长近三个数量级。所以\textbf{在给定尺度与流速下}，磁雷诺数 $R_m=\mu_0Lv/\eta$ 随温度升高而增大：高温聚变等离子体是磁冻结的好近似，低温器件中磁场可以迅速扩散。（注意 $R_m$ 还含 $L$ 与 $v$，扩散时标本身只由 $\eta$ 与尺度决定。）

{\footnotesize
\textbf{出处：}Chen 第 5 章（扩散与电阻率）。
Fitzpatrick 第 2 版第 3、4 章（碰撞与输运）。
}

\textbf{常见错误：}把开尔文直接代入公式（应为 eV）；漏掉 $T^{3/2}$；忽略 $\ln\Lambda$ 的温度依赖（它变化很慢，但并非常数）。
\end{example}


\section{玻姆判据与鞘层厚度}

\begin{example}{离子声速、鞘层厚度与壁面离子通量}{}
\textbf{题目：}氩等离子体 $T_e=3\ \mathrm{eV}$、$n_0=10^{16}\ \mathrm{m^{-3}}$，器壁偏压 $V_w=100\ \mathrm{V}$。
（a）求玻姆速度（离子声速）$c_s$。（b）求\emph{体}等离子体的德拜长度。
（c）用 $s\approx\dfrac{\sqrt2}{3}\lambda_D(n_s)\left(\dfrac{2V_w}{T_e}\right)^{3/4}$ 估算 Child--Langmuir 鞘层厚度（$\lambda_D(n_s)$ 为鞘边密度处的德拜长度）。
（d）求壁面离子通量，并说明它与 (c) 的厚度是否守同一个电流边界条件。

\textbf{解：}
（a）$c_s=\sqrt{\dfrac{k_BT_e}{M_{\mathrm{Ar}}}}
=\sqrt{\dfrac{3\times1.602\times10^{-19}}{6.68\times10^{-26}}}=2.68\times10^{3}\ \mathrm{m\,s^{-1}}$。

（b）\emph{体}等离子体的德拜长度 $\lambda_D(n_0)=7430\sqrt{3/10^{16}}=1.29\times10^{-4}\ \mathrm{m}=0.129\ \mathrm{mm}$。
注意鞘层计算要用\textbf{鞘边}密度 $n_s$ 对应的德拜长度，见 (c)。

（c）鞘层从鞘边密度 $n_s=0.61n_0$ 起步，所以厚度估算式里的 $\lambda_D$ 必须取 $n_s$ 处的值。
$\lambda_D(n_s)=\lambda_D(n_0)/\sqrt{0.61}=1.29\times10^{-4}/0.781=1.65\times10^{-4}\ \mathrm{m}$。
$(2\times100/3)^{0.75}=66.7^{0.75}=23.3$，
$s=0.471\times1.65\times10^{-4}\times23.3=1.81\times10^{-3}\ \mathrm{m}\approx1.8\ \mathrm{mm}$。
注意 $s\gg\lambda_D(n_s)$（约 11 倍）：偏压远大于电子温度时，鞘层远厚于德拜长度。
若误用体密度算出的 $\lambda_D(n_0)$，厚度会被低估 $\sqrt{n_0/n_s}=1.28$ 倍（得到 1.4 mm），并使厚度与 (d) 的通量不守同一个电流边界条件。

（d）鞘边密度 $n_s\approx0.61n_0=6.1\times10^{15}\ \mathrm{m^{-3}}$。
$\Gamma=n_sc_s=6.1\times10^{15}\times2.68\times10^{3}=1.6\times10^{19}\ \mathrm{m^{-2}s^{-1}}$，
即壁面离子电流密度 $j=e\Gamma\approx2.62\ \mathrm{A\,m^{-2}}$。用 (c) 的厚度反代平面无碰撞 Child--Langmuir 关系
\[
j=\frac49\eps_0\sqrt{\frac{2e}{M}}\frac{V_w^{3/2}}{s^2}
\]
得到 $j\approx2.62\ \mathrm{A\,m^{-2}}$（由该式反解鞘厚得 $1.813\ \mathrm{mm}$，与估算式相差 0.2\%）。厚度与通量守同一个电流边界条件。
每个离子在鞘层中被加速获得约 $100\ \mathrm{eV}$ 能量，这正是离子轰击在刻蚀与溅射中起作用的原因。

{\footnotesize
\textbf{出处：}Fitzpatrick 第 2 版第 4 章末 Langmuir Sheaths 与 Langmuir Probes。
Lieberman \& Lichtenberg《Principles of Plasma Discharges and Materials Processing》（按主题：鞘层与玻姆判据）。
}

\textbf{常见错误：}用离子温度计算 $c_s$，而玻姆判据由电子温度决定。把鞘层厚度当成 $\lambda_D$。把\emph{体}密度 $n_0$ 的德拜长度用到鞘边，厚度会被低估 $\sqrt{n_0/n_s}=1.28$ 倍。厚度与壁面通量若用不同密度，两者就不守同一个 Child--Langmuir 电流边界条件。忘记 $n_s\approx0.61n_0$ 这个鞘边密度。
\end{example}


\section{Sweet--Parker 重联速率}

\begin{example}{Sweet--Parker 重联为什么「太慢」}{}
\textbf{题目：}日冕电流片：$T=86\ \mathrm{eV}$（约 $10^{6}\ \mathrm{K}$）、$\ln\Lambda=20$、$L=10^{5}\ \mathrm{m}$、$v_A=10^{5}\ \mathrm{m\,s^{-1}}$。
（a）求电阻率 $\eta$ 与 Lundquist 数 $S=\mu_0Lv_A/\eta$。
（b）求入流速度 $v_{\mathrm{in}}=v_A/\sqrt{S}$ 与重联时间 $\tau_{\mathrm{SP}}=L/v_{\mathrm{in}}$。
（c）与太阳耀斑观测时标（$10^{2}$--$10^{3}\ \mathrm{s}$）比较，能得出什么结论？

\textbf{解：}
（a）$86^{1.5}=7.98\times10^{2}$，
$\eta=5.2\times10^{-5}\times20/7.98\times10^{2}=1.30\times10^{-6}\ \Omega\cdot\mathrm{m}$；
$S=\dfrac{1.2566\times10^{-6}\times10^{5}\times10^{5}}{1.30\times10^{-6}}=9.6\times10^{9}$。

（b）$v_{\mathrm{in}}=v_A/\sqrt{S}=10^{5}/9.82\times10^{4}=1.02\ \mathrm{m\,s^{-1}}$，
$\tau_{\mathrm{SP}}=10^{5}/1.02=9.8\times10^{4}\ \mathrm{s}\approx27\ \mathrm{h}$。

（c）$27$ 小时比观测到的耀斑时标慢约 $10^{2}$--$10^{3}$ 倍。这就是著名的「Sweet--Parker 太慢」问题：全局电阻率太低，重联被限制在极薄的电流片内。解决方向是\textbf{局域化}：撕裂模不稳定性把电流片撕裂成许多小尺度结构（plasmoid 不稳定性），使有效重联率大幅提高。

\textbf{补充关系：}磁场扩散时间 $\tau_\eta=\mu_0L^2/\eta=\sqrt{S}\,\tau_{\mathrm{SP}}=9.6\times10^{9}\ \mathrm{s}$（约 300 年），比 $\tau_{\mathrm{SP}}$ 大 $\sqrt{S}\approx10^{5}$ 倍，这就排除了「整体扩散」的解释，也说明重联必须发生在薄电流片中。

{\footnotesize
\textbf{出处：}Fitzpatrick 第 2 版第 9 章（磁重联：Sweet--Parker 模型、撕裂模、plasmoid 不稳定性）。
}

\textbf{常见错误：}把入流速度取成 $v_A$（$v_A$ 是入流在上游的极限值，重联率取决于 $\sqrt{S}$ 的压制）；漏掉 $S$ 定义中的 $L$ 或 $v_A$；把 $\tau_{\mathrm{SP}}$ 与 $\tau_\eta$ 混为一谈。
\end{example}


\section{八道题与本书章节的对应}

\small
\begin{longtable}{@{}p{1.2cm}p{4.5cm}p{4.3cm}p{4.9cm}@{}}
\caption{附录 K 例题与本书章节、经典教材的对应}\label{tab:app-k-map}\\
\toprule
编号 & 主题 & 本书章节 & 主要出处 \\
\midrule
\endfirsthead
\toprule
编号 & 主题 & 本书章节 & 主要出处 \\
\midrule
\endhead
\midrule
\multicolumn{4}{r}{\small（续下页）}\\
\endfoot
\bottomrule
\endlastfoot
K.1 & 德拜屏蔽与 $N_D$ & 第 2 章 & Chen 第 1 章；Fitzpatrick 第 1 章 \\
K.2 & 磁镜与损失锥 & 第 3 章 & Chen 第 2 章；Fitzpatrick 第 2 章 \\
K.3 & 梯度漂移与曲率漂移 & 第 3 章 & Chen 第 2 章；Fitzpatrick 第 2 章 \\
K.4 & Alfvén 速度与磁冻结 & 第 5 章 & Fitzpatrick 第 8 章；Chen 第 4 章 \\
K.5 & Landau 阻尼 & 第 7 章 & Fitzpatrick 第 7 章；Chen 第 7 章 \\
K.6 & Spitzer 电阻率与扩散 & 第 8 章 & Chen 第 5 章；Fitzpatrick 第 3、4 章 \\
K.7 & 玻姆判据与鞘层 & 第 11 章 & Fitzpatrick 第 4 章；Lieberman \& Lichtenberg \\
K.8 & Sweet--Parker 重联 & 第 9 章 & Fitzpatrick 第 9 章 \\
\end{longtable}
\normalsize

\textbf{做完这八道题之后：}回到各章的自学检查清单（\texttt{checklistbox}）逐条自评，并把做错的题目按附录 I 的三栏法记入错题本。若你发现某道题「看答案全懂、合上书不会」，说明缺的是检索练习而不是理解。按附录 I 的动作五，隔一天再做一遍。
